\documentclass[11pt]{article}
\usepackage[a4paper,margin=1in]{geometry}
\usepackage{amsmath,amssymb,mathtools,bm}
\usepackage{enumitem,booktabs,array}
\usepackage{tcolorbox}
\usepackage{hyperref}
\usepackage[protrusion=true,expansion=false]{microtype}
\hypersetup{colorlinks=true,linkcolor=blue,urlcolor=blue,citecolor=blue}
\setlist[itemize]{topsep=2pt,itemsep=2pt}
\setlist[enumerate]{topsep=2pt,itemsep=2pt}
\newtcolorbox{keybox}{colback=blue!3!white,colframe=blue!40!black,boxrule=0.4pt,arc=1mm}
\newtcolorbox{alertbox}{colback=red!3!white,colframe=red!40!black,boxrule=0.4pt,arc=1mm}
\newtcolorbox{answerbox}{colback=green!3!white,colframe=green!40!black,boxrule=0.4pt,arc=1mm}
\newtcolorbox{drillbox}{colback=yellow!5!white,colframe=orange!60!black,boxrule=0.4pt,arc=1mm}
\newcommand{\EV}{\mathbb{E}}
\newcommand{\Var}{\mathrm{Var}}
\newcommand{\Cov}{\mathrm{Cov}}
\newcommand{\blank}[1]{\underline{\hspace{#1}}}

\title{E05-01: Maksimovic \& Phillips (2001, JF) \\
\large ``The Market for Corporate Assets: Who Engages in Mergers and Asset Sales and Are There Efficiency Gains?'' --- Active Recall Workbook}
\author{Empirical Corporate Finance II --- Exam Prep}
\date{\today}
\begin{document}\maketitle

\begin{keybox}
\textbf{Paper in one sentence.} Using Longitudinal Research Database (LRD) plant-level Census data, MP show that asset sales and partial-firm acquisitions move plants from less-productive to more-productive owners, producing measurable productivity gains post-transaction --- evidence consistent with a neoclassical reallocation theory of M\&A, where the market for corporate assets efficiently allocates capital to its best use.

\textbf{Placement.} Foundational empirical paper for the neoclassical view of M\&A, complementing Harford (2005) on merger waves and motivating productivity-based M\&A studies. Frequently contrasted with managerial-motive explanations (hubris, empire-building).
\end{keybox}

\section*{Round 1: Complete Walk-Through}

\subsection*{1.1 Research question}
Do corporate asset sales and mergers efficiently reallocate assets to more productive owners (``neoclassical'' view) or are they driven by agency/financing frictions or managerial motives? Answer requires plant-level productivity data, not just aggregate firm stock returns.

\subsection*{1.2 Data}
\begin{itemize}
\item Longitudinal Research Database (LRD): plant-level Census of Manufactures + Annual Survey of Manufactures, 1974--1992.
\item $\sim$50,000 plant-year observations linked across owners to track plant transfers.
\item Measure plant-level total factor productivity (TFP) from inputs/outputs.
\item Classify transactions: asset sales (plants transferred outside merger), mergers (firm-level), partial-firm sales.
\end{itemize}

\subsection*{1.3 Empirical design}
\begin{itemize}
\item Compare pre- and post-transaction productivity of transferred plants.
\item Compare buyer and seller average productivity.
\item Condition on industry cycles, economy-wide capital reallocation waves.
\item Match transferred plants to non-transferred controls with similar characteristics.
\end{itemize}

\subsection*{1.4 Headline results}
\begin{itemize}
\item Plants sold are typically less productive than buyer's existing plants in that industry.
\item After transfer, plant productivity increases significantly (+1--3 percentage points TFP).
\item Effect strongest in industries with high productivity dispersion.
\item Buyers are disproportionately main-industry firms; conglomerates often sell plants outside their core.
\item Reallocation is pro-cyclical: more asset transactions in high-growth industries, contra fire-sale theories.
\end{itemize}

\subsection*{1.5 Interpretation}
\begin{itemize}
\item Supports neoclassical reallocation theory: the market for corporate assets moves capital to its highest-value use.
\item Inconsistent with strong ``hubris'' or ``empire-building'' views as dominant motive for mergers.
\item Implies policy interventions restricting asset transactions may reduce efficiency.
\end{itemize}

\subsection*{1.6 Implications}
\begin{itemize}
\item Plant-level productivity is the right unit of analysis for M\&A efficiency.
\item Firm-level event studies miss within-firm reallocation.
\item Industry-cycle variation is crucial: asset sales cluster in times of opportunity, not distress.
\end{itemize}

\subsection*{1.7 Caveats}
\begin{alertbox}
\begin{itemize}
\item LRD coverage limited to manufacturing; service and financial industries excluded.
\item TFP measurement assumptions (Cobb-Douglas, constant returns) may bias results.
\item Selection: plants chosen for sale differ from retained plants on unobservables.
\item Cannot rule out that some transactions are motivated by financial distress or taxes.
\end{itemize}
\end{alertbox}

\section*{Round 2: Level 1 Fill-in}
\begin{enumerate}
\item Data: \blank{2cm} Research Database, plants \blank{1cm}--\blank{1cm}.
\item Productivity measure: plant-level \blank{1cm}.
\item Buyers more productive than \blank{2cm} in target industry.
\item Post-transfer TFP change: $+$\blank{1cm}--\blank{1cm} pp.
\item Theory supported: \blank{2cm} reallocation.
\end{enumerate}

\section*{Round 3: Level 2 Prompts}
\begin{enumerate}
\item Describe the LRD data and why plant-level measurement matters.
\item Report the main productivity-change findings for transferred plants.
\item Explain the neoclassical reallocation theory and how MP test it.
\item Discuss the industry-cycle pattern in asset reallocation.
\item What are the main caveats, and how might they be addressed?
\end{enumerate}

\section*{Round 4: Minimal Scaffolding}
From a blank page: (i) question, (ii) LRD data, (iii) results, (iv) theory, (v) caveats.

\section*{Round 5: Blank Page}
Essay: the market for corporate assets --- evidence for efficient reallocation.

\vspace{6pt}
\begin{drillbox}
\textbf{Speed Drill.} (1) Dataset. (2) Productivity measure. (3) TFP change. (4) Buyer vs.\ seller productivity. (5) Theory supported.
\end{drillbox}

\section*{Quick Reference \& Answer Key}
\begin{answerbox}
\begin{itemize}
\item \textbf{Data.} LRD plant-level panel 1974--92.
\item \textbf{Result.} Plants move less-productive $\to$ more-productive owners.
\item \textbf{Magnitude.} $+1$--$3$ pp TFP post-transfer.
\item \textbf{Theory.} Neoclassical reallocation.
\end{itemize}
\end{answerbox}

\begin{keybox}
\textbf{30-second exam pitch.} Maksimovic \& Phillips (2001) use plant-level Census data (LRD, 1974--1992) to test whether M\&A and asset sales reallocate assets efficiently. Tracking transferred plants, they find that sold plants are typically less productive than the buyer's existing plants in the same industry and, post-transfer, plant TFP rises by $1$--$3$ percentage points, with the strongest effects in industries with high productivity dispersion and during high-growth (not distressed) periods. Buyers concentrate in their main industry; conglomerates tend to shed non-core plants. Results are consistent with a neoclassical reallocation theory of the market for corporate assets --- M\&A moves capital to its highest-value user --- and inconsistent with dominant managerial-motive (hubris, empire-building) interpretations. The paper is a foundational plant-level M\&A study and complements Harford (2005) on merger waves and Rhodes-Kropf et al.\ on valuation-driven M\&A.
\end{keybox}

\end{document}
